\section{Observation Boundaries}\label{sec:observation-boundaries}

The recovery statements concern the public pair \((\alpha,b)\) from
\cref{def:public-observation}. A protocol may instead expose a residue
payload, a digest or a proof about a cochain. Those are different observation
maps. This section records only the mathematical distinctions needed by the
HCP argument. The construction, encoding, visibility policy and security
contracts of the cochain commitment belong to Paper~2
\cite{vanDerVlugt2026protocol}.

\subsection{Exact Cochains and Full Hodge Samples}

Fix one cellular Hodge system and an opening \(a\in C^0\). Its exact
one-cochain \(c=\delta_0a\) satisfies
\[
  P_Hc=P_H\delta_0a=0.
\]
Indeed, \(\image\delta_0\) is orthogonal to the harmonic subspace.
Thus this exact cochain is not the full sample \(b=h+e\) when the harmonic
message \(h\) is nonzero. An encoding modulo an integer, a digest of that
encoding and a proof about its acceptance do not, by their types alone,
supply the full \(b\) or its relation to the operative hidden operator.

There is also a distinction between the exact one-cochain and its opening.
Linearity gives
\[
  \delta_0a=\delta_0a'
  \quad\Longleftrightarrow\quad
  a-a'\in\ker\delta_0.
\]
Consequently the raw coboundary determines an opening only modulo
\(\ker\delta_0\). A unique degree-zero message encoding requires further
application data. This elementary observation is independent of binding a
specified exact \(C^1\) payload cryptographically.

To connect a protocol artifact to the HCP observation, an application must
supply a relation binding that artifact to the full sample, its operative
operator and its classifier value. The artifact can retain an exact
component and separately relate it to \(b\); projecting such a related record
is not an algorithm that extracts the full sample from residue bytes.
Paper~2 gives the concrete payload, snapshot and observation contracts.
The argument here uses only the resulting mathematical pair.

\subsection{The Exact-Component Obstruction}

The loss of information can be quantitative. Under the hypothesis of
\cref{prop:exact-component-obstruction}, a nonzero coexact direction gives
samples with the same exact component and arbitrarily large full Hodge
residual. Hence an observation depending only on that component need not
preserve the full-sample residual or a positive far-radius condition. The
proposition concerns the stated unrestricted Hodge-noise carrier; a
restricted application profile requires its own comparison theorem.

This does not identify the distinct scores
\[
  \|P_{H^\perp}b\|,\qquad
  \operatorname{dist}(b,\image\delta_0),\qquad
  \text{a modular cochain or code distance}.
\]
The first measures disagreement with the harmonic subspace, while the second
measures distance to exactness. The modular distances and their native
normalizations are specified by the protocol profile and analyzed in
Paper~3. No equality of those scores is assumed by the HCP/CVP comparison.

\subsection{Additional Public Observations}

The constructed ambiguity examples fix \((\alpha,b)\). If a larger released
view \(V(x,b)\) is used, the same obstruction applies only after proving
that the two presentations also give equal values of that view. For example,
a view factoring through \((\alpha,b)\) preserves their equality; arbitrary
snapshot roots, full geometric coordinates or proof transcripts need not do
so. No security conclusion for an unspecified larger record follows.

The full projective point of the operative harmonic subspace is a concrete
example of consequential additional information. At the fixed ambient inner
product, it determines that subspace's projector and therefore the HCP target
from \(b\). We prove this reconstruction next. A digest-only view remains a
different observation and is not covered by the reconstruction theorem.
